%! TEX program = xelatex
% ============================================================
% 中文核心期刊论文主文档（manuscript1_journal_article）
% 目标期刊样式参考：《系统工程与电子技术》/《信息与控制》
% 论文核心：多无人机任务分配（task allocation）
% 编译顺序：xelatex -> biber -> xelatex -> xelatex
% 也可直接：latexmk（已提供 .latexmkrc）
% ============================================================
\documentclass[
  zihao   = 5,        % 五号字（期刊常用；如目标期刊要求小四，改 zihao=-4）
  fontset = windows,  % Windows 字体集（宋体/黑体/楷体/仿宋）
                      % Linux/macOS 临时调试可换：fontset=fandol（正式投稿保持 windows）
  scheme  = chinese   % 中文主方案（定义 \songti/\heiti/\zihao 等命令）
]{ctexart}

% 全局宏包与依赖
% ============================================================
% setup/preamble.tex —— 全局宏包与依赖
% 只放"与内容无关"的全局依赖；排版规范见 format.tex
% ============================================================

\usepackage{geometry}          % 页面几何（尺寸在 format.tex 中设置）
\usepackage{graphicx}          % 插图
\usepackage{amsmath,amssymb}   % 数学
\usepackage{amsthm}            % 定理环境（\theoremstyle 依赖，必须加载）
\usepackage{bm}                % 数学粗体
\usepackage{mathtools}         % amsmath 增强
\usepackage{booktabs}          % 三线表
\usepackage{multirow}          % 跨行单元格
\usepackage{caption}           % 图表题定制（格式在 format.tex）
\usepackage{subcaption}        % 子图
\usepackage{xcolor}            % 颜色（占位标记用）
\usepackage{enumitem}          % 列表定制
\usepackage{algorithm}         % 算法浮动体
\usepackage{algpseudocode}     % 算法伪代码

% 参考文献：GB/T 7714—2015 顺序编码制（与 manuscript1 同后端）
\usepackage[
  backend  = biber,
  style    = gb7714-2015,
  sorting  = none              % 按引用顺序编码（顺序编码制）
]{biblatex}
\addbibresource{bib/references.bib}

% hyperref 通常放最后
\usepackage[
  colorlinks = true,
  linkcolor  = black,
  citecolor  = black,
  urlcolor   = blue
]{hyperref}

% 字体（英文 = Times New Roman）
% ============================================================
% setup/fonts.tex —— 字体设置
% 中文：由 ctex 的 fontset=windows 提供（宋体/黑体/楷体/仿宋）
% 英文/数字：Times New Roman（期刊通行要求）
% ============================================================

% 英文主字体 = Times New Roman
% Linux/macOS 无该字体时，临时调试可改用 TeX Gyre Termes（近似 Times）：
%   \setmainfont{TeX Gyre Termes}
\setmainfont{Times New Roman}

% 自定义命令 / 定理环境 / 算法汉化
% ============================================================
% setup/commands.tex —— 自定义命令 / 定理环境 / 算法汉化
% ============================================================

% ---------- 占位标记 ----------
% 正文中所有待写内容统一用 \placeholder{} 标出，终稿前逐项替换并删除。
% 占位文本为灰色斜体，便于一眼识别未完成部分。
\newcommand{\placeholder}[1]{%
  \textcolor{gray}{\textit{【占位：#1】}}}

% 写作要点提示（仅注释用，不输出）
% \writingnote{...} —— 以注释形式写在各节文件里更直观，此命令备用
\newcommand{\writingnote}[1]{}

% ---------- 首页页脚 ----------
% 无编号脚注：用于首页页脚的基金项目 / 通信作者（两刊投稿须知要求）
\newcommand{\blfootnote}[1]{%
  \begingroup
  \renewcommand{\thefootnote}{}%
  \footnote{#1}%
  \addtocounter{footnote}{-1}%
  \endgroup}

% ---------- 常用缩写/排版命令 ----------
\newcommand{\ie}{\textit{i.e.}}
\newcommand{\eg}{\textit{e.g.}}
\newcommand{\etal}{\textit{et al.}}
\newcommand{\R}{\mathbb{R}}          % 实数域
\newcommand{\abs}[1]{\left|#1\right|}
\newcommand{\norm}[1]{\left\|#1\right\|}
\newcommand{\transpose}{\mathrm{T}}      % 转置符号：正体 T（两刊均要求）

% ---------- 定理类环境（汉化） ----------
\theoremstyle{plain}
\newtheorem{definition}{定义}
\newtheorem{theorem}{定理}
\newtheorem{lemma}{引理}
\newtheorem{proposition}{命题}
\newtheorem{assumption}{假设}
\theoremstyle{remark}
\newtheorem{remark}{注}

% ---------- 算法浮动体汉化 ----------
\floatname{algorithm}{算法}
\algrenewcommand{\algorithmicrequire}{\textbf{输入：}}
\algrenewcommand{\algorithmicensure}{\textbf{输出：}}

% 版面格式（页边距 / 标题 / 图表题 / 公式编号）
% ============================================================
% setup/format.tex —— 版面与排版规范（按目标期刊微调）
% 两刊要求对照（详见 README.md）：
%   《系统工程与电子技术》：图表题/注释中文；坐标轴标量和单位；
%     物理量 GB 3100～3102-1993（变量斜体，矢量矩阵黑斜体，
%     集合黑正体，常数/运算符/函数正体）
%   《信息与控制》：图题表题中英双语；公式编号置右；
%     投稿为 Word 双栏（排版终稿由编辑部完成，本模板用于写作）
% 注：各刊细节不同，投稿前以目标期刊"稿件模板/征稿简则"为准。
% ============================================================

% ---------- 页边距（期刊通行：上下 2.5cm，左右 2.0~2.5cm） ----------
\geometry{top=2.5cm, bottom=2.5cm, left=2.2cm, right=2.2cm}

% ---------- 各级标题（黑体；标题层级= 节 / 小节） ----------
\ctexset{
  section = {
    format     = \heiti\zihao{4},      % 节标题：黑体四号
    beforeskip = 1.5ex plus .2ex minus .2ex,
    afterskip  = 1.2ex plus .2ex minus .2ex,
    indent     = 0pt
  },
  subsection = {
    format     = \heiti\zihao{-4},     % 小节标题：黑体小四
    beforeskip = 1.2ex plus .2ex minus .2ex,
    afterskip  = 1ex plus .2ex minus .2ex,
    indent     = 0pt
  },
  subsubsection = {
    format     = \heiti\zihao{5},      % 小小节标题：黑体五号
    beforeskip = 1ex plus .2ex minus .2ex,
    afterskip  = 1ex plus .2ex minus .2ex,
    indent     = 0pt
  }
}

% ---------- 图表题（五号；图题在下、表题在上） ----------
% 《系统工程与电子技术》：图表题和注释用中文
% 《信息与控制》：图题/表题中英双语（\caption{中文题}\caption*{English title}
%   或改用 bicaption 宏包）；投稿时按目标刊切换
\captionsetup{font=small, labelsep=quad, skip=4pt}
\captionsetup[figure]{position=bottom}
\captionsetup[table]{position=top}

% ---------- 公式编号：按节编号 (1)、(2)… ----------
\numberwithin{equation}{section}

% ---------- 段落 ----------
% 中文首行缩进 2 字符（ctex 默认），行距保持紧凑（期刊默认单倍行距）
\linespread{1.1}
\setlength{\parindent}{2em}


% ============================================================
% 论文元信息（投稿前按目标期刊要求逐项填写）
% ============================================================
\newcommand{\paperTitleCN}{【占位：中文标题，一般不超过 20 字】}
\newcommand{\paperTitleEN}{【Placeholder: English Title】}
\newcommand{\paperAuthors}{【占位：作者一$^{1}$，作者二$^{1,2}$】}
\newcommand{\paperAuthorsEN}{【Placeholder: Author One$^{1}$, Author Two$^{1,2}$】}
\newcommand{\paperAffiliation}{%
  （1. 【占位：单位名称】，【占位：城市】 【占位：邮编】；\\
   2. 【占位：单位名称】，【占位：城市】 【占位：邮编】）}
\newcommand{\paperAffiliationEN}{%
  (1. 【Placeholder: Affiliation】, 【City】 【Zip】, China;\\
   2. 【Placeholder: Affiliation】, 【City】 【Zip】, China)}
% 文献标识信息（按目标期刊公布的中图分类号填写）
\newcommand{\clcNumber}{TP391}                 % 中图分类号（占位示例）
\newcommand{\docTypeCode}{A}                   % 文献标志码
\newcommand{\articleNumber}{XXXX-XXXX(YYYY)NN-NNNNN-NN}  % 文章编号（占位）
\newcommand{\doiNumber}{10.xxxxx/xxxxx}        % DOI（占位）

\begin{document}

% ---------- 首页：标题 / 作者 / 单位 / 文献标识 ----------
% ============================================================
% frontmatter/title_info.tex —— 首页信息
% 标题 / 作者 / 单位 / 中图分类号·文献标志码·文章编号·DOI
% 元信息统一在 main.tex 中定义，此处只负责排版。
% ============================================================

\begin{center}
  {\heiti\zihao{3}\paperTitleCN\par}          % 中文标题：黑体三号（≤ 20 字）
  \vspace{1.2em}
  {\kaishu\zihao{-4}\paperAuthors\par}        % 作者：楷体小四
  \vspace{0.8em}
  {\songti\zihao{5}\paperAffiliation\par}     % 单位：宋体五号（省 城市 邮编）
\end{center}

\vspace{1em}

% 文献标识信息行（字号可按期刊要求缩小）
{\songti\zihao{5}\noindent
  中图分类号：\clcNumber\quad
  文献标志码：\docTypeCode\quad
  文章编号：\articleNumber\quad
  DOI：\doiNumber\par}

\vspace{1em}

% ---------- 首页页脚（两刊要求） ----------
% 《系统工程与电子技术》：首页页脚注明基金项目（标准名称+编号）
% 《信息与控制》：首页页脚注明“通信作者：姓名，E-mail”
% \blfootnote 生成无编号首页脚注（定义见 setup/commands.tex）
\blfootnote{基金项目：\fundProject}
\blfootnote{通信作者：\corrAuthor，E-mail：\corrEmail}

% 【注】如目标期刊要求将基金项目/作者简介列于文末，
% 则移除上方脚注，改在 backmatter/author_bio.tex 中列出。

% ---------- 中文摘要 + 关键词 ----------
% ============================================================
% frontmatter/abstract_cn.tex —— 中文摘要 + 关键词
% 撰写要求（两刊对照，投稿时按目标刊取严）：
%   - 《系统工程与电子技术》：150～250 字，关键词 4～8 个
%   - 《信息与控制》：≤ 500 字，关键词 3～6 个（建议用 EI 主题词）
%   - 取严写法：摘要 150～250 字，关键词 4～6 个
%   - 四要素：目的 — 方法 — 结果 — 结论；避免公式/图表/参考文献编号/非常用缩略语
%   - 关键词按“研究对象—方法—性质”顺序排列
% ============================================================

\vspace{1em}

\noindent{\heiti\zihao{5}摘要：}%
    面向任务持续到达、机巢分散部署的广域常态化无人机作业场景，作业过程中常伴随无人机故障、紧急任务等突发扰动。现有动态多无人机任务分配研究多采用事件驱动重分配范式，但普遍隐含“到达任务应尽量全部分配”的假设，较少在机队能力不足时将任务选择作为显式决策；起降结构多简化为单基地或固定返航，终点机巢鲜少进入联合优化。当吞吐受限、机巢分散与地形约束并存时，需要在滚动时域内同时决定执行哪些任务、由何机按何顺序访问以及降落于何巢。
    本文将该问题表述为滚动时域下的选择性多机巢多无人机任务分配：在硬航程与结构约束下优先保障可完成任务的覆盖，覆盖等价时再压缩含返航的航程；允许任务暂缓，并将迟到作为软约束计入代价而非淘汰条件。方法采用两阶段分解。第一阶段在能力上界约束下从动态任务池中确定本周期承诺集，以控制后续组合规模并抑制任务被无限期推迟；第二阶段仅在已选集上求解多机归属、访问顺序与终点机巢，通过搜索与约束修复生成可行航次，使覆盖优先与航程优化在可行域内分层落实。滚动过程吸收新任务到达、平台故障与机巢状态变化，并对受影响子集实施部分重分配，以保持分配方案与动态环境一致。

    （占位）在数字高程模型驱动的多机巢场景及多随机种子实验中，所提框架在可完成任务集上保持高完成率，相对简化距离模型与构造性策略在单位完成航程与时效指标上具有优势，并在故障与紧急任务注入下维持任务守恒与协议一致的重规划行为。结果表明，将任务选择、多机巢终点决策与地形感知代价纳入统一滚动分配框架，可在动态资源约束与突发扰动下实现覆盖优先、成本可控的多机协同调度。

\par\vspace{0.8em}

\noindent{\heiti\zihao{5}关键词：}%
\placeholder{关键词一；关键词二；关键词三；关键词四；关键词五（4～6 个为宜）}

\vspace{1.2em}

% ---------- 英文摘要 + 关键词 ----------
% ============================================================
% frontmatter/abstract_en.tex —— 英文标题 / 作者 / 单位 / Abstract / Keywords
% 撰写要求：
%   - 与中文摘要对应但不是逐字直译，可适度扩充背景
%   - 《系统工程与电子技术》：英文题名实词 ≤ 10 个
%   - 时态：一般现在时 / 现在完成时为主
%   - Keywords 与中文关键词一一对应，全小写（按刊要求）
% ============================================================

\begin{center}
  {\bfseries\zihao{3}\paperTitleEN\par}       % English title
  \vspace{1.2em}
  {\zihao{-4}\paperAuthorsEN\par}             % Authors (pinyin)
  \vspace{0.8em}
  {\zihao{5}\paperAffiliationEN\par}          % Affiliations
\end{center}

\vspace{1em}

\noindent\textbf{Abstract:}%
\placeholder{English abstract. Objective: ...; Methods: ...; Results: ...; Conclusions: ...}

\par\vspace{0.8em}

\noindent\textbf{Keywords:}%
\placeholder{keyword one; keyword two; keyword three; keyword four; keyword five}

\vspace{1.2em}


% ---------- 正文（按节拆分，核心 = 任务分配） ----------
% ============================================================
% sections/s01_introduction.tex —— 1 引言
% 写作要点（四段式，1~1.5 页）：
%   1) 研究背景与意义：多无人机协同应用价值 → 任务分配问题的地位
%   2) 国内外研究现状：按方法脉络综述（精确方法 / 启发式 / 混合……），
%      每类给出代表文献 [x-y]，并指出不足（落脚点=本文要解决的缺口）
%   3) 本文工作：针对上述不足，本文做了什么（3~4 条贡献点）
%   4) 论文结构安排：各节内容一句话概括
% 注意：航迹规划仅作为任务分配的从属/假设出现，不展开综述（论文核心=任务分配）。
% ============================================================

\section{引言}

\placeholder{第一段：研究背景与意义。多无人机（multi-UAV）协同在……领域的应用价值；
任务分配是决定系统效能的关键环节……}

\placeholder{第二段：研究现状综述。按方法类别展开，每段落脚到“现有方法的不足”，
为本文切入点做铺垫。引用文献使用顺序编码 \textbackslash cite\{...\}。}

\placeholder{第三段：本文贡献点（建议用 3~4 条分列或连排）：
① 建立了……任务分配模型；② 提出……求解方法；③ 设计……约束处理机制；
④ 仿真实验验证了……}

\placeholder{第四段：结构安排。第 2 节建立问题模型；第 3 节提出任务分配方法；
第 4 节给出仿真实验与结果分析；第 5 节总结全文。}
          % 1 引言
% ============================================================
% sections/s02_problem_formulation.tex —— 2 问题描述与系统模型
% ============================================================

\section{问题描述与系统模型}

\subsection{问题描述与建模假设}

\placeholder{}

\subsection{符号定义}

\placeholder{}

\subsection{约束条件}

\placeholder{}

\subsection{优化模型与问题性质}

\placeholder{}
   % 2 问题描述与系统模型
% ============================================================
% sections/s03_task_allocation.tex —— 3 滚动时域选择性任务分配方法
% ============================================================

\section{滚动时域选择性任务分配方法}

\subsection{总体框架}

\placeholder{}

\subsection{阶段一：选择性任务接纳机制}

\placeholder{}

\subsection{阶段二：多机巢协同任务分配}


\subsection{滚动闭环与状态继承}

\placeholder{}
       % 3 多无人机任务分配方法（核心章节）
% ============================================================
% sections/s04_experiments.tex —— 4 仿真实验与结果分析
% 写作要点：
%   4.1 实验设置：想定场景、参数取值表、硬件/软件环境、评价指标定义
%   4.2 对比方法与消融设计：对比算法（2~4 个 baseline）+ 本文方法消融变体
%   4.3 结果与分析：主实验表/图 + 定量分析（提升百分比、收敛曲线、
%       方案甘特图/分配结果示意图）
%   4.4 讨论（可选）：敏感性分析（参数影响）、可扩展性（规模增长）
% 注意：实验口径与第 2 节模型约束一一对应；统计显著性/多次独立运行均值±方差。
% ============================================================

\section{仿真实验与结果分析}

\subsection{实验设置}

\placeholder{场景与参数：任务规模、无人机规模、约束参数取值表（表 4.1，三线表）；
运行环境（CPU/内存/语言/求解器）；每组实验独立运行 30 次取均值等口径说明。}

\placeholder{评价指标定义：目标函数值、可行率、求解时间、均衡度等（公式 (4.x)）。}

\subsection{对比方法}

\placeholder{对比算法列表与来源文献（如：遗传算法、差分进化、粒子群、
商业求解器精确解等）；本文方法的消融变体设置（去掉改进点 ①/② 后的版本）。}

\subsection{结果与分析}

\placeholder{主结果表（表 4.2）：各方法在不同规模下的指标对比，
标注最优/次优；定量结论（本文方法较对比方法提升 xx\%）。}

\placeholder{图 4.x：收敛曲线对比；图 4.x：典型场景任务分配结果示意
（无人机-任务连线图或甘特图）。逐图给出分析文字。}

\subsection{敏感性分析}

\placeholder{（可选）关键参数（种群规模、惩罚系数、权重系数等）
对结果的影响曲线 + 结论。}
           % 4 仿真实验与结果分析
% ============================================================
% sections/s05_conclusion.tex —— 5 结论
% 写作要点（半页以内）：
%   1) 工作总结：3~4 条，与引言贡献点呼应（用过去时陈述做了什么）
%   2) 主要结论：定量/定性结论各 1~2 条
%   3) 展望：1~2 句，不铺开（如动态任务分配、航迹-任务联合优化等）
% 注意：不引入新内容、不重复摘要原文。
% ============================================================

\section{结论}

\placeholder{本文针对……问题，建立了……模型，提出了……方法，
设计了……机制。仿真实验表明：① ……；② ……；③ ……。}

\placeholder{下一步工作：……（一两句展望，如动态不确定环境下的在线任务分配、
任务分配与航迹规划联合优化等）。}
            % 5 结论

% ---------- 参考文献（GB/T 7714—2015，顺序编码制） ----------
\printbibliography[title={参考文献}]

% ---------- 收稿日期 / 基金项目 / 作者简介（部分期刊置于文末） ----------
% ============================================================
% backmatter/author_bio.tex —— 文末信息
% 收稿日期 / 修回日期 / 基金项目 / 作者简介
% 注：《系统工程与电子技术》《信息与控制》等多将此块置于文末；
% 如目标期刊要求首页脚注，请将本块移至 frontmatter/title_info.tex。
% ============================================================

\vspace{1.5em}
\hrule
\vspace{0.8em}

\noindent{\heiti\zihao{5}收稿日期：}\placeholder{YYYY-MM-DD}
\quad{\heiti\zihao{5}修回日期：}\placeholder{YYYY-MM-DD}

\par\vspace{0.8em}

\noindent{\heiti\zihao{5}基金项目：}%
\placeholder{国家自然科学基金（编号：XXXXXXX）；XX 省重点研发计划（编号：XXXXXXX）等}

\par\vspace{0.8em}

\noindent{\heiti\zihao{5}作者简介：}%
\placeholder{姓名（出生年—），性别，籍贯，职称/学位，研究方向。
示例：张三（1998—），男，湖北武汉人，硕士研究生，研究方向为多无人机协同任务分配。}

\par\vspace{0.5em}

\noindent{\heiti\zihao{5}通信作者：}%
\placeholder{姓名，E-mail：xxx@xxx.edu.cn}


\end{document}
